\mxunit{automatismes-guide}{Guide d'utilisation (version professeur)}
\subsection*{Principe}
Chaque séance de 2G ouvre sur une fiche de cinq questions (5 à 8 minutes), alignée sur la séance du calendrier du jour :
\begin{enumerate}[label=\textbf{\arabic*.},leftmargin=1.6em]
  \item \textbf{Échauffement} : un prérequis de cycle~4 directement utile à la séance.
  \item \textbf{Cœur} : une question sur la notion du jour (ou de la veille).
  \item \textbf{Cœur +} : une seconde question plus exigeante (QCM, vrai/faux justifié, lecture de code, méthode).
  \item \textbf{Retour} : une question hors chapitre, issue d'un attendu « à consolider » ou « à entretenir » (espacement des révisions).
  \item \textbf{Fil rouge} : une question ciblée sur un attendu en « renforcement prioritaire » du test de positionnement du 29/09/2026.
\end{enumerate}
En version professeur, chaque question est précédée de son étiquette de diagnostic (par exemple \mxtag{EQ}) ; la réponse attendue figure en rouge, suivie d'une remarque ou d'un piège fréquent.

\subsection*{Déroulement suggéré (7 min)}
\begin{enumerate}[label=\textbf{\arabic*.},leftmargin=1.6em]
  \item 3 min : réponses individuelles à l'ardoise ou sur le cahier de rituels, sans calculatrice sauf mention.
  \item 2 min : mise en commun ; faire verbaliser une réponse \emph{fausse} et une réponse \emph{juste} (« comment as-tu fait\,? »).
  \item 2 min : correction ciblée des questions 4--5 (celles du diagnostic) avec la phrase-méthode associée.
\end{enumerate}

\subsection*{Variante FLE/FLS}
\begin{itemize}[leftmargin=1.2em,itemsep=2pt]
  \item \textbf{Consigne} : projeter la question avec le verbe d'action surligné (\emph{calculer, résoudre, montrer, comparer}) et un exemple de réponse modèle la première semaine.
  \item \textbf{Phrases-cadres} : « Je calcule\ldots{} donc\ldots{} », « Je compare\ldots{} car\ldots{} », « Le contre-exemple est\ldots{} ». Les afficher en bandeau permanent.
  \item \textbf{Glossaire visuel} : fiche mots-clés (déterminant, colinéaire, médiane, taux d'évolution, fréquence conditionnelle) avec schéma et traduction dans la langue d'origine de l'élève si utile.
  \item \textbf{Binômes} : 30 secondes de reformulation entre pairs avant de répondre ; l'élève allophone peut d'abord répondre par un geste ou un schéma, puis formaliser à l'oral.
  \item \textbf{Écrits} : accepter les notations symboliques avant la phrase rédigée, puis exiger la phrase de conclusion une fois sur deux.
\end{itemize}

\subsection*{Grille d'analyse des erreurs}
{\footnotesize\mxarraystretch{1.2}
\begin{tabular}{@{}p{2.2cm} p{4.4cm} p{4.8cm} p{4.4cm}@{}}\toprule
Étiquettes & Erreur observée & Indice dans la réponse & Geste de remédiation \\ \midrule
EQ, CAL & Erreur de signe en transposant ou en distribuant & $5x-7=2x+8\Rightarrow3x=1$ & Écrire chaque étape ; vérifier en remplaçant. \\
DEV, LIT & Distributivité oubliée ($2(x+3)=2x+3$) & Résultat non réduit ou faux pour $x=1$ & Test numérique ; aire d'un rectangle. \\
PUI, SUB & Confusion $(-a)^2$ et $-a^2$ & Réponse négative pour un carré & Parenthèses systématiques à la substitution. \\
GEO, REP & Distance sans racine ; $x_A-x_B$ au lieu de $x_B-x_A$ & $AB=25$ ou signe incohérent & Vérifier l'ordre de grandeur sur un schéma. \\
PRO & Addition de taux successifs ; confusion proportion/évolution & $+20\,\%$ puis $-10\,\%=+10\,\%$ & Passer par les coefficients multiplicateurs. \\
PRB, LEC & Confusion $P_A(B)$ et $P_B(A)$ & Mauvais total de référence & Entourer le total du sous-ensemble « sachant que ». \\
REG, IMG & Confusion image et antécédent ; lecture inversée & Réponse sur le mauvais axe & Colorier : abscisse = départ, ordonnée = arrivée. \\
DEM, LIT & Un exemple pris pour une preuve & « Ça marche pour $n=3$ » & Distinguer test, contre-exemple et démonstration. \\
ALG & Exécution de code sans suivre les variables & Valeur finale incorrecte & Tableau de suivi (variable / valeur / étape). \\ \bottomrule
\end{tabular}}

\subsection*{Barème}
Les rituels sont formatifs : aucun barème n'est joint. Un barème détaillé par compétence (Chercher, Modéliser, Représenter, Raisonner, Calculer, Communiquer) peut être ajouté si l'une de ces fiches sert de support d'évaluation.

\subsection*{Sources et points de vigilance}
\begin{itemize}[leftmargin=1.2em,itemsep=2pt]
  \item Programme de mathématiques de seconde générale et technologique, arrêté du 26-2-2026 (NOR MENE2602914A), BO n°14 du 2 avril 2026 : rubrique « Automatismes » et parties « Nombres et calculs », « Géométrie », « Fonctions », « Statistiques et probabilités », « Vocabulaire ensembliste et logique ».
  \item Calendrier : consulté le 06/10/2026 ; l'ordre des notions suit le calendrier et non les dates du document de progression, qui diffèrent (ex.\ fonctions : 3--15 déc.\ 2026 au calendrier).
  \item Exercices : \textbf{créés pour ce document} avec votre autorisation ; toutes les valeurs numériques et algébriques ont été recalculées par script (voir \texttt{tools/}), les raisonnements ont été relus. Les données du tableau « précipitations » (fiche~35) sont fictives.
  \item Convention des quartiles : plus petite valeur dont au moins $25\,\%$ (resp.\ $75\,\%$) des données lui sont inférieures ou égales (rang $\lceil n/4\rceil$, $\lceil3n/4\rceil$) --- \textbf{[VERIFY]} coordonner avec les conventions du manuel utilisé.
  \item Les questions sur les solides, patrons et perspective (étiquettes SOL, PER) relèvent du cycle~4 : le programme de Seconde ne comporte pas de chapitre dédié ; elles servent le diagnostic.
  \item Fiche~77 (question~3) : notion de taux d'accroissement en anticipation de la Première, hors chapitre de Seconde.
  \item Liste des automatismes évaluables à l'épreuve anticipée (BO n°24 du 12 juin 2025, session 2027) : \textbf{[VERIFY --- check Eduscol]} pour la session 2028.
\end{itemize}
